\section{Overview}
\label{sec:overview}

We introduce HADL through two self-contained workflows. The first covers types and policies, showing how values, types, and authorization rules produced by the LLM are validated under the same type discipline and authorization semantics that govern the rest of the program. The second covers workflow generation, showing how a \texttt{gen} call can produce a complete HADL workflow that executes as part of the workflow that produced it. Together, the two workflows cover HADL's main contributions before the formal treatment in Sections~\ref{sec:design} and~\ref{sec:implementation}.

HADL programs consist of workflow declarations, which are typed functions that compose four domain-specific primitives with standard imperative constructs. The primitive \texttt{ask} prompts a human operator and returns a \texttt{String}. The primitive \texttt{say} prints a string to the output stream and returns unit. The primitive \texttt{gen} sends a description to an LLM gen session configured to return structured data and requires an explicit type annotation at the binding site. The primitive \texttt{agent} sends an instruction to an LLM reasoning session and returns a \texttt{String}. These four primitives are first-class expressions in the language, not callbacks or graph nodes, so they compose naturally with \texttt{let} and \texttt{var} bindings, \texttt{if}/\texttt{else} branches, \texttt{while} loops, \texttt{for} iteration, and inline JavaScript expressions via the \texttt{js} keyword. The type checker runs before execution begins, rejecting \texttt{gen} calls without type annotations, assignments to immutable variables, and return type mismatches, all with diagnostic source spans.

\subsection{Gradual Validation for Types and Policies}
\label{sec:overview-types-policies}

Many agentic workflows need to handle two kinds of runtime information at once. The shape of the data flowing between \texttt{gen} calls is not known until the LLM produces it, and the authorization rules that govern those calls are often specific to the operator, the task, or the stage of the workflow. HADL handles both with the same type discipline. The \texttt{Schema} type is compatible with every static type and resolves to a concrete type once the LLM delivers one, after which the rest of the program is checked under the same type system. The \texttt{Policy} type binds an LLM-generated Cedar policy as a value, and an explicit \texttt{enforce} expression activates it, after which every subsequent \texttt{gen} or \texttt{agent} invocation is checked against a principal supplied at the call site.

Figure~\ref{fig:clinical-trial} shows a clinical trial workflow. The \texttt{@retries 3} pragma on line~1 activates self-healing, described below. Line~4 prompts the operator for a trial description. Line~5 binds \texttt{crf} with a \texttt{Schema} annotation, signalling to the type checker that this variable will hold a type produced at runtime. When line~5 executes, the interpreter parses the LLM's response as a JSON Schema value, converts it to a concrete HADL type, and checks the remainder of the program against that concrete type under the same type discipline that governs the rest of the language. Line~6 produces a blinding policy and line~8 activates it with \texttt{enforce}. After line~8, every \texttt{gen} and \texttt{agent} invocation constructs an authorization triplet with its declared principal and is checked against the installed policy before any LLM call is made. Line~10 asks the LLM for a list of patient identifiers, bound as \texttt{Schema} since the array structure is determined at runtime. The \texttt{for} loop on lines~12--16 iterates that list. Line~13 binds \texttt{visit} with type \texttt{crf}, so the output is validated against the concrete schema produced on line~5, and the field access \texttt{visit.patient\_id} on line~15 is checked under the same type system. The \texttt{gen} call inside the loop carries the \texttt{as role} clause, so each invocation is authorized for the principal the operator declared on line~9, and line~14 counts each successful visit in the mutable \texttt{total} via an inline \texttt{js} expression. If a serious adverse event is reported, lines~20--21 produce and activate a second policy. The two policies are merged append-only, so access rules only tighten. Line~22 can only run when the principal is \texttt{safety\_board}, because the merged policy blocks every other role from adverse event reports.

The \texttt{@retries} pragma activates self-healing. When validation rejects an LLM output, the runtime identifies the \texttt{gen} call that produced the rejected value, re-invokes it with the violated constraint appended to the original prompt, and re-runs the check under the same type and authorization discipline. The loop runs up to the declared retry limit. If the retry limit is exhausted, the interpreter raises the underlying error with the full diagnostic.

\subsection{Gradual Validation for Workflow Generation}
\label{sec:overview-workflow}

Agentic workflows often delegate subtasks whose structure is itself the subject of the current reasoning step. A scoring function, a triage procedure, or a data cleaning pipeline may all be worth producing on the fly, shaped by what the operator has just said. HADL handles this by making workflows first-class values with arrow types. When a \texttt{gen} call is annotated with an arrow type, the LLM produces a complete HADL workflow as a JSON-serialized AST, and the \texttt{eval} expression deserializes and invokes it in the calling workflow's environment. The generated workflow is not a separate program. It runs as part of the workflow that produced it, sharing the caller's variable scope, and every \texttt{eval} site is checked against the arrow's parameter and return types.

Figure~\ref{fig:overview-workflow} shows a car recommendation workflow and the workflow it produces. Line~3 of the outer workflow generates a \texttt{Schema} describing a single car. Line~4 asks the LLM to generate a scorer with the arrow type \texttt{(car\_schema -> Number)}, an arrow whose parameter type is the runtime-generated \texttt{car\_schema} itself. The LLM produces the workflow shown below the arrow, taking a car conforming to the schema and returning a numeric score. The \texttt{while} loop on lines~9--18 fills the schema three times, binding a fresh \texttt{car} value with type \texttt{car\_schema} on each iteration so that every LLM output is validated against the concrete shape produced on line~3. Line~11 applies the generated scorer via \texttt{eval scorer(car)}. At the \texttt{eval} expression the type checker verifies that the argument type matches the arrow's parameter type and that the result is used where the return type is accepted. Line~14 extracts the car's name via an inline \texttt{js} expression when a new best score is found. The arrow between the two listings indicates that the generated workflow is produced by the \texttt{gen} call on line~4 of the outer workflow and runs in the same interpreter context. The generated workflow itself calls \texttt{gen} three times to rate comfort, cost, and reliability, and uses an inline \texttt{js} expression to average the three ratings.

\begin{figure}[H]
\centering
\begin{subfigure}{\linewidth}
\begin{lstlisting}[language=hadl,showstringspaces=false]
workflow pick_car() -> String {
  let prefs = ask "What matters most?"
  let car_schema: Schema = gen "Schema for a car with name and type"
  let scorer: (car_schema -> Number) = gen "Workflow that
    scores a car on `prefs`"
  var best: String = ""
  var top: Number = 0
  var i: Number = 0
  while (i < 3) {
    let car: car_schema = gen "Fill the schema with a popular car"
    let score: Number = eval scorer(car)
    if (score > top) {
      top = score
      best = js car.name
    }
    say "`car.name` rated `score`"
    i = js i + 1
  }
  say "Recommended `best`"
  best
}
\end{lstlisting}
\subcaption{Outer workflow.}
\end{subfigure}

\vspace{0.4em}
\begin{tikzpicture}[>=Stealth]
  \draw[->, very thick] (0,0.3) -- (0,-0.3);
  \node[right, font=\footnotesize] at (0.15,0)
    {generated at runtime, runs in the outer workflow's scope};
\end{tikzpicture}
\vspace{0.4em}

\begin{subfigure}{\linewidth}
\begin{lstlisting}[language=hadl,showstringspaces=false]
// produced by the gen call on line 4 above
workflow score(car: car_schema) -> Number {
  let comfort: Number = gen "Rate comfort of `car.name`"
  let cost: Number = gen "Rate cost of `car.name`"
  let reliability: Number = gen "Rate reliability of `car.name`"
  js (comfort + cost + reliability) / 3
}
\end{lstlisting}
\subcaption{Generated workflow.}
\end{subfigure}
\caption{A car recommendation workflow illustrating gradual validation for workflow generation. Line~3 of the outer workflow generates a schema for a single car. Line~4 asks the LLM for a scorer with arrow type \texttt{(car\_schema -> Number)}, the parameter being the runtime-generated type itself. The LLM produces the workflow shown below, taking a car that conforms to the schema and returning a numeric score. The \texttt{while} loop fills the schema three times, and \texttt{eval scorer(car)} on line~11 invokes the generated workflow for each fresh car, evaluating it in the outer workflow's variable scope. At every \texttt{eval} site, the type checker verifies that the argument matches the arrow's parameter type and that the result is used where the return type is accepted. The generated workflow itself calls \texttt{gen} three times and averages the results with an inline \texttt{js} expression, showing that a generated workflow has full access to HADL's primitives and JavaScript interop.}
\label{fig:overview-workflow}
\end{figure}
